Validation loss became the default checkpoint-selection rule in robot learning by inheritance,
not by argument. It works in supervised learning because a classifier is graded one example at
a time, exactly what the loss measures. A policy is graded by a journey, and that is where the
inheritance breaks: per-step errors the loss forgives compound into trajectories that leave the
training manifold and fail. We have named this the validation gap and given it a calibration
treatment borrowed from surrogate-endpoint validation, with a pre-registered protocol that
fixes the tasks, two architectures, eight offline signals, rollout count, and analysis before
any data is seen.

The contribution does not depend on which way the data falls. If a coherence-based offline
signal recovers the deployment ranking, training pipelines can drop a costly rollout sweep. If
none does past a horizon threshold, the field has a calibrated reason to budget for rollouts
exactly where it matters. Either way, practitioners gain a rule rather than a folk belief, and
the released benchmark lets anyone test a new offline signal against ground-truth success
without retraining. The broader lesson generalises beyond this study: stop grading sequential
decision-makers one step at a time, and measure the property that deployment actually rewards.
